% ============================================================
% [SUSPECT FLAG (2026-04-23 設置) → CLEAR (2026-05-01 03:10 JST)]
% ============================================================
% Generated within 4/18-4/22 thinking-extended degrade window (K9_118.7).
% Probe matrix verification (B2 designer, 2026-05-01) indirect-cleared:
% K9_118 batch content-correct via §S-SPL exact match + §S-CB-diag
% superior-to-o3 quality. Significance statement is 150w summary-style
% (minimal reasoning depth requirement) and inherits batch clearance.
% Reference: 6回目作業用_出力はここだけにしろ/notes/note_wave61_suspect_flag_disposition_2026-05-01.md
% Reference: alert_20260501_0310_JST_..._wave61_suspect_flag_clear_content_diff_verdict_SA_unfreeze.md
% SA 投稿可 (SA submission unfreeze, 2026-05-01 03:10 JST).
% V2 = V1 + this CLEAR header update (no body change).
% ============================================================
%
% Paper 4 Significance Statement V1 --- Science Advances submission
% Derived from K9_118.7 (wave60 response), Q2 V1 150w SA tier (selection matrix score 9.5/10).
% 2026-04-23 wave61 integration by Ａ１ (加藤方程式９). SUSPECT (see flag above).
%
% Source basis: paper4_main_V51.tex.txt, paper4_SI_V65.tex.txt, paper4_refs_V23.bib.txt
% (confirmed integer-resolution breakdown 17/22, composition 21/22, continuous H_eff 22/22,
% patent H=2, GDP H=4-5 corridor, MAE=0.0308, RMSE=0.0447).
%
% Primary version: V1 (150 words, Science Advances tier, default for SA submission).
% Alternative escalation versions (V2/V3/V4) preserved as comment block below
% for PNAS resubmission (V2, 120 words strict), Nature escalation (V3, 97 words, lay-facing),
% and Science escalation (V4, 135 words, editor-facing impact paragraph).

\documentclass[11pt]{article}
\usepackage[margin=1in]{geometry}
\usepackage{amsmath,amssymb}
\usepackage[T1]{fontenc}
\usepackage[utf8]{inputenc}
\setlength{\parindent}{0pt}
\setlength{\parskip}{0.45em}
\pagestyle{empty}

\begin{document}

\section*{Significance Statement V1 --- Science Advances (150 words)}

Cities concentrate people, infrastructure, and ideas, but urban scaling remains hard to use predictively unless the exponents themselves can be explained. On the canonical Bettencourt 22-indicator benchmark, integer relay-depth predictions account for 17 reported confidence intervals but leave five structured misses, all in crossover regions rather than scattered noise. Paper~4 introduces the Kato equation, a zero-parameter exponent ladder that maps an integer network relay depth to mirror-symmetric superlinear and sublinear predictions. Mixed observables follow a share-weighted composition rule; composition lifts confidence-interval coverage to 21/22, and continuous \(H_{\mathrm{eff}}\) resolves the final miss without indicator-specific slope fitting. Patent scaling remains an \(H=2\) innovation case, GDP and wages occupy an \(H=4\)--\(5\) corridor, and infrastructure follows the negative branch. The result is a network-anchored scaling framework that organizes superlinear production, linear balances, and sublinear infrastructure within one falsifiable urban-physics scaffold for comparison across cities, sectors, definitions, and measurement frames globally.

\vspace{1em}
\noindent\textbf{Word count:} 150 \hfill \textbf{Target tier:} Science Advances

\end{document}

% ============================================================
% ALTERNATIVE ESCALATION VERSIONS (do NOT submit with V1)
% Source: K9_118.7 Q3/Q4/Q5. Kept for PNAS/Nature/Science escalation handoff.
% ============================================================
%
% ---------- V2 (PNAS resubmission, strict 120 words) ----------
% Cities are where population, infrastructure, and ideas concentrate, yet
% urban science still lacks a simple reason why different city indicators
% scale with different exponents. Paper 4 proposes that each exponent is
% selected by network relay depth: how many serial layers an observable
% must pass through before it becomes measured output. The resulting Kato
% ladder predicts superlinear, linear, and sublinear cases with one rule.
% On the canonical 22-indicator Bettencourt benchmark, integer relay-depth
% predictions match 17 confidence intervals; the five misses are structured
% crossover cases, not random failures. A composition rule raises coverage
% to 21/22, and continuous H_eff resolves the remaining miss without tuning
% slopes for individual indicators. The framework turns urban scaling from
% fitted curves into a testable network-based theory.
% [Word count: 120]
%
% ---------- V3 (Nature escalation, lay-facing, 97 words) ----------
% Cities shape jobs, roads, energy use, and innovation, but planners often
% see urban scaling as a set of fitted curves rather than a guide to why
% cities behave differently. Paper 4 offers a simple interpretation: an
% indicator scales according to the number of network steps needed to
% produce it. Innovation, wages, housing, roads, and utilities therefore
% fall on related rungs of one ladder instead of separate laws. That
% matters because policy can ask which network depth an intervention
% changes, not only whether a city is large. The framework could make urban
% planning more predictive, comparable, and testable.
% [Word count: 97]
%
% ---------- V4 (Science escalation, editor paragraph, 135 words) ----------
% Urban scaling has carried a 20-year puzzle: why do patents, wages, GDP,
% housing, roads, cables, and utilities obey stable power laws but land on
% different exponents? Paper 4 resolves this by replacing fitted slopes
% with a zero-parameter relay-depth ladder. The Kato equation assigns each
% observable an integer H determined by network relay depth, then predicts
% beta_+(H) for production, beta_-(H) for infrastructure, and balanced
% mixtures for near-linear quantities. On Bettencourt's 22-indicator
% benchmark, integer assignments recover 17 confidence intervals;
% composition raises the score to 21/22, and continuous H_eff explains the
% final crossover case without indicator-specific fitting. The same
% framework unifies superlinear, linear, and sublinear indicators, places
% GDP and wages in a deep H=4-5 corridor, and keeps patents at H=2. Its key
% claim is experimentally sharp: measure topology first, then predict the
% exponent.
% [Word count: 135]
%
% ============================================================
% SELECTION MATRIX (K9_118.7 Q6)
% ============================================================
%   Science Advances submission / editor pitch  --> V1 (9.5/10) [THIS FILE]
%   PNAS resubmission                           --> V2 (9/10)
%   Nature escalation                           --> V3 (8.5/10)
%   Science escalation                          --> V4 (9/10)
%   General cover-letter insert                 --> V1 or V4 (9/10)
% ============================================================
% Common core, one sentence:
%   Urban scaling exponents are not arbitrary fitted slopes; they are
%   selected by network relay depth through a zero-parameter Kato ladder,
%   with composition and continuous H_eff explaining crossover indicators
%   across the Bettencourt benchmark.
% ============================================================
